\documentclass[10pt,a4paper]{amsart}
\usepackage[margin=19mm]{geometry}
\usepackage{amsmath,amssymb,mathtools}
\usepackage[scr=rsfs,cal=euler]{mathalfa}
\usepackage{xfrac}
\usepackage[unicode,hidelinks,bookmarks=false]{hyperref}
\newcommand{\ie}{i.e.\ }
\newcommand{\cf}{cf.\ }
\newcommand{\resp}{respectively\ }
\newcommand{\Subsection}{\subsection}
\providecommand{\nobreakdash}{\nobreak\hbox{-}}
\setlength{\emergencystretch}{2em}
\multlinegap0pt
\usepackage{tikz}
\usepackage{amssymb}
\newtheorem{corollary}{Corollary}
\title{Boundary value problems for linear differential-algebraic equations: solvability via the Kronecker canonical form}
\author{Anar Assanova, Carsten Trunk, Roza Uteshova, Henrik Winkler}
\date{Source: arXiv:2609.07980v1 (CC BY 4.0)}
\begin{document}
\maketitle

\noindent\emph{Source and licence.} Boundary value problems for linear differential-algebraic equations: solvability via the Kronecker canonical form. Version arXiv:2609.07980v1 (CC BY 4.0), distributed by arXiv under the Creative Commons Attribution 4.0 International licence, http://creativecommons.org/licenses/by/4.0/, as recorded in the arXiv licence metadata for this exact version and shown on the abstract page https://arxiv.org/abs/2609.07980v1. Full licence notice: \texttt{licenses/arxiv-2609.07980-CC-BY-4.0.txt}.\par\smallskip

\noindent\textit{Editorial scope.} Counted unit: the article's corollary corollary\_ODEpart\_uniqueness (source0.tex lines 489--498). The article's complete original proof of it is delivered with it (source0.tex lines 499--506, one \texttt{proof} environment of 8 lines). No mathematics is written, repaired or re-derived: the statement and the proof are the article's own lines, in the article's own notation, and no retained line is substituted or re-flowed.  Retained uncounted as the source-local context the proof needs: lines 349--351; lines 420--420; lines 421--422; lines 433--435; lines 438--440; lines 477--479.  Nothing was omitted from any retained span, so the package carries no omission record.  No retained line contains a citation, so the wrapper carries no bibliography and no bibliography entry.  No retained line contains a figure environment, an image-inclusion command or drawing code, so no image is delivered and the package declares no asset dependency.  Editorial formatting only, in addition: an AMS \texttt{amsart} preamble that restates the article's own declarations and macros used by the retained lines, namely the environments \texttt{corollary} on separate counters, together with the standard packages those lines need.  The retained environments print in this document as Corollary 1 in order of appearance, whereas the article prints the same items with the numbering of its own full document.  The complete native source of the article as distributed in the arXiv e-print for this exact version is bundled unchanged as \texttt{sources/209682/source0.tex}.
\begin{equation}\label{Alina}
A_L^{-1} :=(A^\top A)^{-1}A^\top
\end{equation}

\begin{equation}\label{x_difeq}\dot{x}_0(t)=A_0x_0(t)+f_0(t), \end{equation}

   \begin{equation}\label{bc_dif}
    B_0x_0(0)+C_0x_0(T)=d_0,\end{equation}

has a solution $x_0\in C^1([0,T],\mathbb{R}^{n_0})$ if and only if \begin{equation}\label{d_0_orth}
    \widehat{d}_0\in \textnormal{ran}~ Q_0.
    \end{equation}

\begin{equation}\label{sol_x0_gen}
    x_0(t)=\int\limits_0^t e^{A_0(t-s)}f_0(s)ds+e^{A_0t}\left[Q_0^g\widehat{d}_0+(I_{n_0}-Q_0^g Q_0)z\right],
\end{equation} where  $z$ is an arbitrary vector in $\mathbb{R}^{n_0}$.

\begin{equation}\label{Q_syst_difpart}
    Q_0\lambda=\widehat{d}_0.
\end{equation}

\begin{corollary}\label{corollary_ODEpart_uniqueness}
  Let $f_0\in C([0,T],\mathbb{R}^{n_0})$. Then problem \eqref{x_difeq}, \eqref{bc_dif} is uniquely solvable if and only if, in addition to \eqref{d_0_orth},
\begin{equation}\label{rankQ_0}
    \emph{rank}~Q_0=n_0,
\end{equation}
in which case the solution is given by
\begin{equation}\label{sol_x0_unique}
    x_0(t)=\int\limits_0^t e^{A_0(t-s)}f_0(s)ds+e^{A_0t}(Q_0^\top Q_0)^{-1}Q_0^\top\widehat{d}_0.
\end{equation}
\end{corollary}

\begin{proof}
If condition \eqref{rankQ_0} is satisfied, then the matrix $Q_0$ has full column rank, implying that its kernel contains only the zero vector. Consequently, the system of equations  \eqref{Q_syst_difpart} can have at most one solution. Moreover, the fulfillment of condition \eqref{d_0_orth} ensures the existence of a solution.

Since $Q_0$ has full column rank, it admits a left inverse, denoted by $(Q_0)_L^{-1}$, which serves as a generalized inverse of  $Q_0$.
The unique solution is given by expression \eqref{sol_x0_unique}, which is derived from \eqref{sol_x0_gen} by substituting $Q_0^g$ with
$(Q_0)_L^{-1}=(Q_0^\top Q_0)^{-1}Q_0^\top$, see
\eqref{Alina}.
\end{proof}

\end{document}
